\chapter{Basket and ETF Arbitrage}
\label{ch:basket_etf_arb}
\chaptermark{Basket and ETF arbitrage}

An exchange-traded fund (ETF) on the S\&P~$500$ holds $500$
underlying stocks in declared proportions. Summing those prices at
the declared weights gives the \emph{synthetic} value of one ETF
share: what the wrapper should cost if it equalled its contents.
The ETF's market price can drift from synthetic value, but not far:
\emph{authorised participants} (a class of large dealer) can
mechanically convert constituents into ETF shares (and back)
through the issuer, pocketing the discrepancy. This chapter is
about what happens when that mechanism is, and is not, available.

In Prosperity~$3$ the basket trade was the largest single source of
PnL for many top teams. \texttt{PICNIC\_BASKET1} contained six
\texttt{CROISSANTS}, three \texttt{JAMS}, and one \texttt{DJEMBE}
per share; \texttt{PICNIC\_BASKET2} contained four
\texttt{CROISSANTS} and two \texttt{JAMS}. The key caveat: the
Prosperity simulator exposes \emph{no} creation / redemption
mechanism. PICNIC\_BASKET prices mean-revert toward synthetic value
only because the simulator's data-generating process is built that
way, not because anyone can force convergence. The Prosperity
basket trade is a \emph{pure cointegration play with a known hedge
ratio} (the cleanest application of
\cref{ch:pairs_cointegration}), but not, strictly, an arbitrage.

The chapter has two halves. The first develops the theory: creation
/ redemption, the synthetic-value formula, static-vs-dynamic hedge
ratios, and Avellaneda and Lee's $2010$ PCA-residual framework.
The second half is a case study of the Frankfurt Hedgehogs'
$2$nd-place PICNIC\_BASKET strategy from Prosperity~$3$. Every
numerical claim traces to
\texttt{imc-prosperity-3/README.md} or
\texttt{imc-prosperity-3/FrankfurtHedgehogs\_polished.py}.

\begin{backgroundbox}[Prerequisites]
From earlier:
\begin{itemize}
  \item \Cref{ch:mean_reversion_ou}: OU $z$-score machinery,
  half-life $h_{1/2} = (\ln 2)/\theta$, stationary std
  $\sigma_\infty = \volat/\sqrt{2\theta}$, entry/exit/stop-out
  rule, AR(1) estimator. The basket-basis rule is the same
  $z$-score applied to the residual.
  \item \Cref{ch:pairs_cointegration}: cointegration and
  cointegration vectors, Engle--Granger two-step, and
  \emph{declared} (given by the basket constructor) vs
  \emph{estimated} (must fit and walk-forward refit) hedge ratios.
  \item \Cref{ch:price_value}: wall-mid as the reference for
  Prosperity spreads; the \cref{ch:participants} taxonomy for who
  the APs of a real ETF are.
  \item Forward: \cref{ch:kelly_risk} for the
  variance-vs-expected-return reasoning behind Round~$5$
  half-hedging; \cref{ch:ml_trading} for the ML view of the
  Avellaneda--Lee residual.
\end{itemize}
From outside the book:
\begin{itemize}
  \item Linear combinations of random variables: if
  $B = \sum_i w_i C_i + \epsilon$ then
  $\Var[B] = \sum_i w_i^2 \Var[C_i] + \Var[\epsilon] +
  2 \sum_{i<j} w_i w_j \Cov[C_i, C_j]$; a portfolio $+1$ basket,
  $-w_i$ of each constituent has variance $\Var[\epsilon]$ alone.
  \item Fund-structure vocabulary: a fund is a legal vehicle, the
  issuer runs it, its ``shares'' trade on an exchange. ETFs
  differ from mutual funds in trading intraday rather than only
  at end-of-day net asset value (NAV) through the issuer.
\end{itemize}
\end{backgroundbox}

\section{ETFs and the creation/redemption mechanism}
\label{sec:ch11_creation_redemption}

An ETF alone is not an arbitrage instrument; it is a
passively-managed fund whose shares trade on an exchange. The
arbitrage angle comes from a parallel mechanism in the prospectus
that lets a specific class of dealer convert wrapper to contents
and back at near-zero cost.

\begin{definition}[Exchange-traded fund]
\index{exchange-traded fund (ETF)|textbf}\index{ETF|textbf}
\label{def:etf}
An \emph{exchange-traded fund (ETF)} is a passively-managed
investment fund whose shares trade intraday on a securities
exchange. The fund holds a basket of underlying assets (stocks,
bonds, commodities, futures) in declared proportions, the
\emph{creation basket}, and the published net asset value (NAV)
per ETF share equals the sum of the constituents' market prices
weighted by those proportions, less management fees. SPY (the
SPDR~S\&P~$500$ ETF), QQQ (Invesco Nasdaq-$100$), and IWM
(iShares Russell~$2000$) are the three most-traded examples in
US equities.
\end{definition}

\begin{definition}[Authorised participant]
\index{authorised participant (AP)|textbf}\index{AP (authorised participant)|textbf}
\label{def:ap}
An \emph{authorised participant (AP)} is a large broker-dealer
contractually permitted by the ETF issuer to deliver the creation
basket of constituents to the fund and receive new ETF shares in
return (\emph{creation}), or to deliver ETF shares back to the
fund and receive the constituents (\emph{redemption}). APs
typically pay a small fixed fee per creation / redemption unit
(thousands of US dollars on a US-equity ETF, irrespective of
notional). For SPY the AP is one of a handful of US bank
broker-dealers; the contract is a private bilateral agreement
between AP and issuer. A retail trader cannot create or redeem:
the barrier is both legal (no AP contract) and operational
(delivering a $500$-name basket in the correct proportions, to
the issuer's custodian, on the right settlement cycle, requires
the clearing and prime-brokerage plumbing only large dealers
have).
\end{definition}

The mechanism gives the AP an automatic arbitrage. When the ETF
trades \emph{above} synthetic value the wrapper is said to be at
a \emph{premium}; \emph{below} synthetic value, at a
\emph{discount}. A premium typically arises when end-investor
demand for ETF shares outruns the pace at which APs are creating
new units; the arbitrage below is precisely the mechanism that
collapses it. Suppose the ETF trades at $P^{\text{ETF}}_t$ with
synthetic value
$P^{\text{synth}}_t = \sum_i w_i C_{i,t} < P^{\text{ETF}}_t$
(a premium). The AP can:

\begin{enumerate}
  \item Buy the constituents at cost $P^{\text{synth}}_t$ per
  creation unit (plus bid--ask and impact, small for liquid names).
  \item Deliver one creation unit to the fund, plus the creation fee.
  \item Receive newly-issued ETF shares.
  \item Sell them on the exchange at $P^{\text{ETF}}_t$.
\end{enumerate}

Profit per creation unit is
$P^{\text{ETF}}_t - P^{\text{synth}}_t - (\text{fee} + \text{costs})$.
Step~$4$'s selling pressure pushes $P^{\text{ETF}}_t$ back toward
$P^{\text{synth}}_t$; the incentive vanishes once the gap falls
below the fee-plus-cost floor. Redemption is the mirror image: when
$P^{\text{ETF}}_t < P^{\text{synth}}_t$ the AP buys cheap ETF
shares, hands them to the fund, receives constituents, and sells
them at the higher synthetic value.

\begin{figure}[ht]
\centering
\begin{tikzpicture}[
  box/.style={draw, rounded corners=4pt, minimum width=3.2cm,
    minimum height=1cm, text centered, font=\small, align=center},
  arr/.style={-Stealth, thick},
]
\node[box, fill=blue!15]   (ap)   at (0,0)  {Authorised\\Participant (AP)};
\node[box, fill=green!15]  (fund) at (7,0)  {ETF Fund\\(issuer)};
\node[box, fill=gray!18]   (mkt)  at (0,-3) {Stock\\Market};
\node[box, fill=orange!15] (exch) at (7,-3) {Exchange\\(ETF shares)};
% Creation (top, going right)
\draw[arr, blue!65]
  ([yshift= 7pt]ap.east) -- node[above, font=\footnotesize]{deliver basket}
  ([yshift= 7pt]fund.west);
\draw[arr, green!65]
  ([yshift=-7pt]fund.west) -- node[below, font=\footnotesize]{receive new ETF shares}
  ([yshift=-7pt]ap.east);
% AP sources basket from market
\draw[arr, gray!60] (mkt.north) -- node[left, align=center, font=\footnotesize]{buys\\constituents} (ap.south);
% AP sells ETF on exchange
\draw[arr, dashed, orange!70]
  (ap.south east) -- node[below, sloped, font=\footnotesize]{sells ETF at premium}
  (exch.north west);
% Fund arrow to exchange (conceptual)
\draw[arr, dashed, gray!50] (fund.south) -- node[right, align=center, font=\footnotesize]{shares\\trade here} (exch.north);
% Label
\node[font=\footnotesize\itshape, text=blue!60] at (3.5, 0.7) {Creation flow};
\end{tikzpicture}
\caption{ETF creation mechanism. When the ETF trades at a \emph{premium} to
NAV, the authorised participant buys the constituent basket on the stock market,
delivers it to the fund, and receives freshly minted ETF shares to sell on
the exchange, locking in the premium as riskless profit.
Redemption is the reverse: the AP buys ETF shares cheaply, redeems them for
the basket, and sells the constituents.
This in-kind arbitrage keeps the ETF price anchored to its net asset value.}
\label{fig:etf_creation}
\end{figure}

\begin{keyfactbox}[The ETF basis is bounded by the round-trip creation/redemption fee]
Define the \emph{ETF basis} as
\begin{equation}
  \text{Basis}_t \;\defeq\; P^{\text{ETF}}_t - P^{\text{synth}}_t,
  \qquad P^{\text{synth}}_t \defeq \sum_{i=1}^n w_i\,C_{i,t}.
  \label{eq:etf_basis}
\end{equation}
On a real, liquid ETF, the absolute basis is bounded above by
the round-trip cost an AP would pay to close the gap:
\begin{equation}
  |\text{Basis}_t|
  \;\lesssim\;
  \underbrace{\text{creation fee}}_{\text{fixed per unit}}
  \;+\;
  \underbrace{\text{constituent half-spread}}_{\text{liquidity cost}}
  \;+\;
  \underbrace{\text{market impact}}_{\text{size-dependent}}.
  \label{eq:basis_bound}
\end{equation}
SPY and QQQ in normal markets: $1$--$5$ bps. Illiquid sector or
emerging-market ETFs whose constituents trade in a different time
zone: tens of bps or more.
\end{keyfactbox}

\begin{intuitionbox}
Creation / redemption is a \emph{hard floor} on what would
otherwise be a soft mean-reversion trade. The OU residual of
\cref{ch:mean_reversion_ou} was bounded only by the stationary
distribution: the spread could wander out arbitrarily long before
reverting, stopped only by statistics and a stop-out. The
cointegration spread of \cref{ch:pairs_cointegration} has the same
problem; the LTCM box there is the canonical example.

The ETF basis on a liquid AP-served ETF is different. Once it
exceeds the fee-plus-cost floor, APs converge within minutes
(often seconds); the gap closes as mechanical engineering, not
statistical drift. ETF basis trading is among the safest
stat-arb variants and, as the case study shows, one of the
lowest-margin, because the floor is well-known and well-policed.
\end{intuitionbox}

\section{The basket spread formula}
\label{sec:ch11_spread_formula}

Strip away the wrapper and the trade is a linear combination of
prices. Let $B$ be the basket, $C_1, \ldots, C_n$ its constituents,
with declared weights $w_1, \ldots, w_n$ (non-negative integers in
Prosperity; fractional on a real ETF).

\begin{definition}[Synthetic basket value and basket basis]
\index{synthetic basket value|textbf}\index{basket basis|textbf}
\label{def:basket_basis}
The \emph{synthetic value} of one basket share at time $t$ is
\begin{equation}
  B^{\text{synth}}_t \;\defeq\; \sum_{i=1}^n w_i\,C_{i,t},
  \label{eq:synthetic}
\end{equation}
and the \emph{basket basis} (or simply \emph{basket spread}) is
\begin{equation}
  \text{Basis}_t \;\defeq\; B_t - B^{\text{synth}}_t
  \;=\; B_t - \sum_{i=1}^n w_i\,C_{i,t}.
  \label{eq:basis}
\end{equation}
The vector
$\bm\beta = (1,\,-w_1,\,-w_2,\,\ldots,\,-w_n)\in\R^{n+1}$ is the
\emph{declared cointegration vector}: the $\R^{n+1}$ portfolio
$+1$ basket, $-w_i$ of each constituent, has time-$t$ value
exactly $\text{Basis}_t$.
\end{definition}

\begin{remark}
The basket's cointegration vector is \emph{declared by
construction}: the issuer publishes the weights, so there is no
Engle--Granger or Johansen fit. The ADF test on the residual stays
informative (it signals whether wrapper and contents have
decoupled), but the slope is given.
\end{remark}

Under cointegration (enforced by creation / redemption on a real
ETF; by the simulator's data-generating process in Prosperity) the
basis is stationary around some mean, often but not always zero, so
the OU $z$-score machinery of \cref{ch:mean_reversion_ou} applies
directly. Compute sample mean $\hat\mu_{\text{Basis}}$ and standard
deviation $\hat\sigma_{\text{Basis}}$, define the running $z$-score
\begin{equation}
  z_t \;\defeq\;
  \frac{\text{Basis}_t - \hat\mu_{\text{Basis}}}
       {\hat\sigma_{\text{Basis}}},
  \label{eq:basis_zscore}
\end{equation}
and apply the threshold rule:

\begin{keyfactbox}[Basket basis trading rule]
With thresholds $z_{\text{entry}} > z_{\text{exit}} \geq 0$, apply
the \cref{ch:mean_reversion_ou} rule:
\begin{itemize}
  \item $z_t > +z_{\text{entry}}$: \emph{short the basis}, i.e.\ short
  $1$ basket, long $w_i$ of each constituent. Basket is rich.
  \item $z_t < -z_{\text{entry}}$: \emph{long the basis}, i.e.\ long
  $1$ basket, short $w_i$ of each constituent. Basket is cheap.
  \item $|z_t| < z_{\text{exit}}$: close; the basis has reverted.
\end{itemize}
The hedge leg is sized \emph{at the declared weights} because
$\bm\beta$ is declared: no slope refit, only the running mean and
variance of the residual need updating.
\end{keyfactbox}

Expected round-trip duration and PnL follow
\cref{ch:mean_reversion_ou} applied to the basis residual: the
basket spread strategy is the OU $z$-score strategy with a
multi-asset hedge leg.

\subsection{A toy numerical example}
\label{sec:ch11_toy_example}

Take a basket $B$ with two constituents, weights $w_1 = 3$,
$w_2 = 2$, and wall-mid prices at three consecutive timesteps:
\begin{center}
\begin{tabular}{rccc}
\toprule
$t$ & $B_t$ & $C_{1,t}$ & $C_{2,t}$ \\
\midrule
$1$ & $510$ & $100$ & $105$ \\
$2$ & $508$ & $101$ & $103$ \\
$3$ & $516$ & $100$ & $106$ \\
\bottomrule
\end{tabular}
\end{center}

With $B^{\text{synth}}_t = 3 C_{1,t} + 2 C_{2,t}$ and
$\text{Basis}_t = B_t - B^{\text{synth}}_t$:
\begin{align*}
  B^{\text{synth}}_1 &= 3 \cdot 100 + 2 \cdot 105 = 510, &
  \text{Basis}_1 &= 510 - 510 = 0, \\
  B^{\text{synth}}_2 &= 3 \cdot 101 + 2 \cdot 103 = 509, &
  \text{Basis}_2 &= 508 - 509 = -1, \\
  B^{\text{synth}}_3 &= 3 \cdot 100 + 2 \cdot 106 = 512, &
  \text{Basis}_3 &= 516 - 512 = +4.
\end{align*}
The sample mean of $\{0, -1, +4\}$ is $\hat\mu = 1$ and the sample
standard deviation (unbiased $n - 1$ divisor) is
$\hat\sigma = \sqrt{((0-1)^2 + (-1-1)^2 + (4-1)^2)/2}
= \sqrt{7} \approx 2.65$. The $z$-score at $t = 3$ is
\begin{equation*}
  z_3 \;=\; \frac{4 - 1}{2.65} \;\approx\; 1.13.
\end{equation*}
With $z_{\text{entry}} = 1.0$ this crosses and the rule fires
\emph{short basis}: sell one $B$, buy three $C_1$ and two $C_2$.
If the basis reverts to its mean of $1$ next bar (constituents
unchanged), PnL is $(\text{Basis}_3 - \text{Basis}_4) \cdot 1 = 3$
seashells per basket short, hedge leg breaking even.

Two takeaways. First, the \emph{hedge ratio is the literal share
count}: no fitting. Second, $\hat\mu = 1 \neq 0$; the $z$-score
is taken against it, not against zero (the running-mean correction
of \cref{sec:ch11_hh_premium_drift}). Shorting whenever
$\text{Basis}_t > 0$ would have shorted at $t = 1$ at zero, sat
through $-1$, and unwound at $+4$ for a four-seashell loss per
unit.

\section{Static vs dynamic hedge ratios}
\label{sec:ch11_static_dynamic}

One degree of freedom hides in \eqref{eq:synthetic}: how do you
know the weights $w_i$? Real ETF: the issuer publishes, the AP
trades. Synthetic basket without a mechanical wrapper-contents
link (proprietary index, factor-mimicking portfolio): murkier.

\begin{definition}[Static hedge ratio]
\index{hedge ratio!static|textbf}
\label{def:static_hedge}
A \emph{static hedge ratio} fixes $w_i$ once at the constructor's
declared values (or, on a real ETF, the prospectus) and never
re-estimates.
\end{definition}

\begin{definition}[Dynamic hedge ratio]
\index{hedge ratio!dynamic|textbf}
\label{def:dynamic_hedge}
A \emph{dynamic hedge ratio} estimates $w_i$ empirically via
Engle--Granger or Johansen (\cref{ch:pairs_cointegration}) on a
rolling window, refit on a fixed cadence ($W = 250$, $R = 21$
trading-day defaults). The hedge ratio at time $t$ is the latest
fit, held out of sample until the next refit.
\end{definition}

The right choice depends on whether wrapper and contents are
mechanically coupled.

\paragraph{Static hedge: when the wrapper is the contents.} On a
real ETF the AP loop closes the basis at the published weights;
any departure leaves exposure the AP will not bail out. Same for
Prosperity PICNIC\_BASKETs: weights $(6, 3, 1)$ and $(4, 2, 0)$
are declared by the simulator; any re-estimate adds noise.

\paragraph{Dynamic hedge: when the wrapper is loosely coupled.} A
proprietary ``smart-beta'' factor index with no AP, rebalancing
at month-end, can exhibit a drifting empirical hedge ratio even
with published methodology. A dynamic hedge on rolling windows
tracks the drift; a static hedge anchored to the month-end
snapshot accumulates unhedged factor exposure the trader
mistakes for alpha.

\begin{pitfallbox}[Static hedge is right only when you can actually do creation/redemption]
The temptation is to plug any basket's published weights into
\eqref{eq:synthetic} and call it done. Correct for SPY, QQQ, and
Prosperity PICNIC\_BASKETs, where creation/redemption or the
simulator enforces the weights. \emph{Wrong} for a non-tradable
proprietary index whose weights drift, for an actively-managed
mutual fund with monthly disclosures, or for any basket whose
``published weights'' are marketing, not a tradable contract.

Diagnostic: \emph{can someone, somewhere, do creation / redemption
against this basket at the published weights?} If yes, static
hedge. If no, fit with Engle--Granger or Johansen
(\cref{ch:pairs_cointegration}) and walk-forward refit. The
half-second check catches this trap before shipping a strategy
silently long an unintended factor.
\end{pitfallbox}

\section{Case study: Frankfurt Hedgehogs and PICNIC\_BASKET}
\label{sec:ch11_hedgehogs}

The Frankfurt Hedgehogs placed $2$nd in IMC Prosperity~$3$ (2025);
their PICNIC\_BASKET strategy
\citep[\texttt{imc-prosperity-3/README.md}, \S\,Round 2]{frankfurthedgehogs2025}
was the leaderboard's largest basket-PnL source. Their public
source \texttt{FrankfurtHedgehogs\_polished.py} lets us cross-check
every parameter; every numerical claim below traces to the writeup
or to a source line.

\subsection{The setup}
\label{sec:ch11_hh_setup}

Round~$2$ introduced three individual products
(\texttt{CROISSANTS}, \texttt{JAMS}, \texttt{DJEMBES}) and two
baskets:

\begin{itemize}
  \item \texttt{PICNIC\_BASKET1}: $6$ \texttt{CROISSANTS}~$+$~%
        $3$ \texttt{JAMS}~$+$~$1$ \texttt{DJEMBE}.
  \item \texttt{PICNIC\_BASKET2}: $4$ \texttt{CROISSANTS}~$+$~%
        $2$ \texttt{JAMS}.
\end{itemize}

\texttt{DJEMBES} appears only in basket~$1$, which matters for the
inter-basket spread below. Position limits (\texttt{POS\_LIMITS}
dict, source lines $30$--$44$): $60$ basket~$1$, $100$ basket~$2$,
$250$ \texttt{CROISSANTS}, $350$ \texttt{JAMS}, $60$
\texttt{DJEMBES}.

\begin{pitfallbox}[Prosperity has no creation / redemption mechanism]
\label{pitfall:no_create_redeem}
The PICNIC\_BASKET in Prosperity~$3$ cannot be unwrapped into
constituents: no AP, no issuer, no creation fee. The basket is a
tradable instrument the simulator's bots quote; its price tracks
synthetic value only because the data-generating process adds a
mean-reverting noise term on top of the constituent-implied price
(Frankfurt writeup: ``\,\emph{the most natural generation process
seemed to be that the three constituents' prices were
independently randomized, and a mean-reverting noise sequence was
then added on top to produce the basket price}\,'').

The basis is therefore bounded only by the \emph{statistical}
mean reversion of that noise term, not by the hard floor
\eqref{eq:basis_bound} sets on a real ETF. The Hedgehogs' strategy
is a pure cointegration play (\cref{ch:pairs_cointegration}), not
a hard arbitrage (\cref{ch:arbitrage_zoo}); every LTCM-box risk
lesson carries over.
\end{pitfallbox}

\subsection{Spread construction}
\label{sec:ch11_hh_spread_construction}

The Hedgehogs computed two separate spreads, one per basket. With
$P^{B1}_t$ and $P^{B2}_t$ the basket prices and
$P^{\text{cr}}_t$, $P^{\text{jam}}_t$, $P^{\text{dj}}_t$ the
constituent prices, define
\begin{align}
  S^{B1}_t &\defeq P^{B1}_t \;-\; \bigl(6\,P^{\text{cr}}_t \;+\; 3\,P^{\text{jam}}_t \;+\; 1\,P^{\text{dj}}_t\bigr),
  \label{eq:s_b1}\\
  S^{B2}_t &\defeq P^{B2}_t \;-\; \bigl(4\,P^{\text{cr}}_t \;+\; 2\,P^{\text{jam}}_t\bigr).
  \label{eq:s_b2}
\end{align}
This is the static-hedge basket basis of
\cref{def:basket_basis} with simulator-declared weights. The
source hard-codes them as
\texttt{ETF\_CONSTITUENT\_FACTORS = [[6, 3, 1], [4, 2, 0]]} (line
$57$); line $437$ computes the spread as a numpy dot of
constituent wall-mids with the factor vector, subtracted from the
basket wall-mid (line $440$). The wall-mid is the
\cref{ch:price_value} wall-mid (midpoint of deepest visible bid
and ask), not top-of-book mid.

The Hedgehogs also considered an \emph{inter-basket} spread
\begin{equation}
  S^{\text{inter}}_t
  \;\defeq\; P^{B1}_t \;-\; 1.5\,P^{B2}_t \;-\; P^{\text{dj}}_t,
  \label{eq:s_inter}
\end{equation}
which removes the constituents by combining the two baskets.
Since $1.5 \times (4, 2, 0) = (6, 3, 0)$, the difference
$(6, 3, 1) - 1.5 \times (4, 2, 0) = (0, 0, 1)$ leaves one
\texttt{DJEMBE} unhedged, hence the $-P^{\text{dj}}_t$. The
writeup is explicit (\S\,Round 2): ``\emph{we quickly identified
that comparing baskets directly to their constituent sums was the
stronger and more reliable path}.'' Production trades only
$S^{B1}_t$ and $S^{B2}_t$. Likely reason: constituent legs are
deep while \texttt{PICNIC\_BASKET2} has thin top-of-book volume.

\subsection{Empirical spread behaviour}
\label{sec:ch11_hh_empirical}

Figure~$4$ (writeup, \S\,Round 2) plots $S^{B1}_t$ and $S^{B2}_t$
across a round. Two dominant features:

\begin{enumerate}
  \item Short-term mean reversion at a ``noise around a level''
  frequency, consistent with an OU residual on top of a stationary
  mean.
  \item A persistent positive premium: $S^{B1}_t$'s running mean
  is non-zero. The wrapper trades above the sum of its parts by
  tens of seashells, the level drifting slowly across a round.
\end{enumerate}

The source initialises premiums as
\texttt{INITIAL\_ETF\_PREMIUMS = [5, 53]} (line $62$, a prior from
Round~$1$ back-test medians), then updates online
(\cref{sec:ch11_hh_threshold}).

The writeup publishes no values for OU half-life or stationary
std; Figure~$4$ suggests a half-life of tens of timesteps and a
stationary std on the order of the premium itself.

\subsection{The static threshold rule}
\label{sec:ch11_hh_threshold}

The Hedgehogs applied a fixed-threshold rule. With base thresholds
$T^{B1}, T^{B2} > 0$ and running premium $\hat\mu^{Bk}_t$, the
de-meaned spread is
\begin{equation}
  \tilde S^{Bk}_t \;\defeq\; S^{Bk}_t - \hat\mu^{Bk}_t,
  \label{eq:demeaned_spread}
\end{equation}
and:

\begin{itemize}
  \item $\tilde S^{Bk}_t > +T^{Bk}$: short the basket (sell bid
  wall, position-limit sized);
  \item $\tilde S^{Bk}_t < -T^{Bk}$: long the basket (buy ask
  wall, position-limit sized);
  \item $\tilde S^{Bk}_t$ crosses zero through an open position:
  close.
\end{itemize}

The zero-crossing close is gated by
\texttt{ETF\_CLOSE\_AT\_ZERO = True} (line $67$), the closing
branch being the ``\texttt{elif ETF\_CLOSE\_AT\_ZERO}'' block on
lines $493$--$501$. Base thresholds are
\texttt{BASKET\_THRESHOLDS = [80, 50]} (line $59$): enter
basket~$1$ at $\pm 80$ from the running mean, basket~$2$ at
$\pm 50$. These are \emph{absolute-seashell} thresholds, not $z$-scores.
The Hedgehogs chose raw thresholds because basis volatility did
not appear to vary meaningfully: ``\emph{using a $z$-score
normalizes the spread by recent volatility, but unless volatility
is known to vary meaningfully over time \ldots\ this introduces
unnecessary complexity and risk of overfitting}'' (writeup,
\S\,Round 2).

\paragraph{Why $80$ and $50$, not $200$ and $100$?} A small grid
search (Figure~$5$, the ``Optimal Parameter Search Grid'' heat map)
led the Hedgehogs to \emph{not pick the best back-test
parameter} but the centre of a flat plateau on the PnL surface,
the canonical defence against small-sample overfitting. The
chrispyroberts P3 writeup\footnote{See
\texttt{research/writeups/chrispyroberts\_P3/ROUND 2/basket\_gridsearch\_results.csv};
an independent grid search confirms a broad optimum, not a sharp
peak.} reaches a similar conclusion.

\subsection{Dynamic threshold adjustment using Olivia}
\label{sec:ch11_hh_olivia}

The strategy's most distinctive feature, the one that
materially separated its PnL from a vanilla threshold rule, is
the use of an external trading-bot signal to bias the thresholds.
The bot is \emph{Olivia}, already detected on \texttt{SQUID\_INK}
in Round~$1$ as an ``informed-buyer / informed-seller'' whose
trades preceded short-term moves. In Round~$2$ the same pattern
applied to \texttt{CROISSANTS}: Olivia long predicted a bullish
\texttt{CROISSANTS} move and hence, since \texttt{CROISSANTS} is
the largest constituent of both baskets, on the basis.

The source mechanic (lines $480$--$491$) is a threshold
\emph{shift}, not a $z$-score adjustment. With
$d_t \in \{-1, 0, +1\}$ for Olivia's direction and magnitude $A$
(\texttt{ETF\_THR\_INFORMED\_ADJS = [90, 90]} line $65$), the
modified entry rule is
\begin{equation}
  \text{enter short basket if}\quad
    \tilde S^{Bk}_t \;>\; +T^{Bk} + d_t \cdot A,
  \label{eq:olivia_short}
\end{equation}
\begin{equation}
  \text{enter long basket if}\quad
    \tilde S^{Bk}_t \;<\; -T^{Bk} + d_t \cdot A.
  \label{eq:olivia_long}
\end{equation}
The single offset $d_t \cdot A$ shifts \emph{both} thresholds the
same way. Olivia long ($d_t = +1$) slides the band up: more
positive deviation for a short, less negative for a long, so the
strategy biases toward long-basket. Olivia short ($d_t = -1$)
biases toward short-basket. $d_t = 0$ reduces to
\cref{sec:ch11_hh_threshold}.

\begin{prosperitybox}[title={Prosperity example: the Olivia cross-product adjustment}]
\label{prosp:olivia}
The clearest illustration is in the Frankfurt writeup itself
(\S\,Round 2): \emph{``if our base threshold was $\pm 50$,
detecting Olivia as short would shift the long entry to $-80$
and the short entry to $+20$.''} Plug $T = 50$, $d_t = -1$,
$A = 30$ into \eqref{eq:olivia_short}--\eqref{eq:olivia_long}:
\begin{align*}
  \text{long entry}  &= -T + d_t A \;=\; -50 + (-1)(30) = -80,\\
  \text{short entry} &= +T + d_t A \;=\; +50 + (-1)(30) = +20.
\end{align*}
Matches the writeup. An Olivia long symmetrically shifts long
entry to $-20$, short entry to $+80$. Production $A = 90$ is
larger than the didactic $30$ because production base thresholds
are $80$ and $50$ (not $50$ and $50$); $A$ was re-tuned on the
larger back-test grid.

The economic content is cross-product information leakage.
\texttt{CROISSANTS} is the dominant constituent in both baskets
($6$ in B$1$, $4$ in B$2$), so any informed signal on
\texttt{CROISSANTS} translates almost mechanically into one on the
basis. Shifting the threshold band rather than the running-mean
estimate keeps the two estimators separate: $\hat\mu^{Bk}_t$
updates from price data alone
(\cref{sec:ch11_hh_premium_drift}); offset $d_t \cdot A$ comes
from Olivia alone; they combine at execution.

The contribution is large. The Hedgehogs report (writeup,
\S\,Round 2 conclusion): ``\emph{about $40{,}000$--$60{,}000$
SeaShells per round on baskets, plus another $20{,}000$ SeaShells
per round directly from trading Croissants individually}.''
Attribution in \cref{sec:ch11_hh_pnl}.
\end{prosperitybox}

\subsection{The drifting premium and the running-mean correction}
\label{sec:ch11_hh_premium_drift}

$S^{B1}_t$'s running mean was non-zero: a persistent positive
premium drifting slowly across a round. A zero-centred rule is
biased: the basis spends more time positive, the strategy runs
net short on average, long-entry never fires. Fix: estimate
$\hat\mu^{Bk}_t$ online and subtract before thresholding.

The estimator (line $448$) is an \emph{incremental running mean}:
\begin{equation}
  \hat\mu^{Bk}_n
  \;=\; \hat\mu^{Bk}_{n-1}
  \;+\; \frac{S^{Bk}_n - \hat\mu^{Bk}_{n-1}}{n}
  \label{eq:running_mean}
\end{equation}
with $n$ incrementing from \texttt{n\_hist\_samples = 60{,}000}
(line $61$) and $\hat\mu^{Bk}_0 =$
\texttt{INITIAL\_ETF\_PREMIUMS = [5, 53]} (line $62$). Not an EMA
(the most-recent-observation weight is $1/n$, not a fixed
$\alpha$); the $60{,}000$-sample prior dominates, so live updates
move the estimate slowly.

\texttt{CALCULATE\_RUNNING\_ETF\_PREMIUM = True} (line $68$)
gates the live update; False keeps only the prior. Production ran
with it on.

\begin{pitfallbox}[The basket premium is not constant]
\label{pitfall:basket_premium}
The most expensive beginner failure: compute spread as
$B - \sum_i w_i C_i$, centre the threshold at zero, trade. If the
basket carries a persistent premium (Prosperity's
PICNIC\_BASKET$1$ did, by tens of seashells), the spread is
almost always positive and only the ``short basket'' branch
fires. The trader becomes systematically short, correct in
expectation \emph{only if} the premium is about to revert to zero,
which it is not.

Fix: the running-mean update \eqref{eq:running_mean}, subtracted
\emph{before} thresholding. \cref{ch:mean_reversion_ou}'s
$z$-score does this implicitly via $\hat\mu$ subtraction; raw
thresholds make the step explicit.

Deeper lesson: a basket basis with non-zero mean can still be
cointegrated. The residual's mean encodes fees, microstructure
costs, and adverse-selection premia, and is non-zero in general.
Centre on the running mean, not structural zero.
\end{pitfallbox}

\subsection{PnL attribution}
\label{sec:ch11_hh_pnl}

The headline PnL is from the writeup's Round~$2$ conclusion:
``\emph{about $40{,}000$--$60{,}000$ SeaShells per round on
baskets, plus another $20{,}000$ SeaShells per round directly
from trading Croissants individually}.'' We treat it as ground
truth: it is the only number published, and it cannot be
re-attributed without rerunning the simulator.

Qualitative observations:

\begin{itemize}
  \item Basket-leg to \texttt{CROISSANTS}-only PnL ratio $\approx
  2$--$3$:$1$. Basket dominates; Olivia on \texttt{CROISSANTS}
  alone is still a meaningful book.
  \item The $40$k--$60$k range spans the five rounds, re-tuned
  each round.
  \item PnL is seashells per $10{,}000$-timestep round; position
  limits above.
\end{itemize}

Economic anchor: with stationary basis range $\pm O(100)$ around
the premium, capturing a fraction of each round-trip over several
dozen cycles yields $\sim 10^4$ seashells per round per basket.
Olivia contributes the \emph{bias} that tilts each round-trip
toward the predicted side.

\subsection{The Round 5 half-hedging decision}
\label{sec:ch11_hh_half_hedge}

In the final round the Hedgehogs had a
$\sim$$190{,}000$-seashell lead; full-aggression upside was small
versus the downside of a spread excursion closing the gap. Their
decision (writeup, \S\,Round 5): \emph{half-hedge}. Buy
$\tfrac{1}{2} w_i$ of each constituent against each long basket
(sell $\tfrac{1}{2} w_i$ against each short).

Source: \texttt{ETF\_HEDGE\_FACTOR = 0.5} (line $70$), fed into
line $532$ as
\texttt{expected\_hedge\_position += -basket.expected\_position
* etf\_const\_factor * ETF\_HEDGE\_FACTOR}.

The reasoning is the linear-combination identity of
\cref{ch:kelly_risk}. Let $X$ be unhedged PnL (basis signal, no
hedge) and $Y$ fully-hedged PnL ($w_i$ of each constituent). Any
convex combination $\alpha X + (1-\alpha) Y$ is a valid strategy.
The Hedgehogs chose $\alpha = \tfrac{1}{2}$. Then:

\begin{itemize}
  \item \textbf{Expected return} is linear:
  $\E[\alpha X + (1 - \alpha) Y] = \alpha \E[X] + (1-\alpha) \E[Y]$.
  Both books trade the same signal with positive expected PnL, so
  the half-hedged book splits the difference. In particular,
  $\E[Y] > 0$: hedging removes constituent-direction exposure;
  the basis, the alpha, is preserved.
  \item \textbf{Variance} is not linear:
  \begin{equation}
    \Var[\alpha X + (1-\alpha) Y]
    = \alpha^2 \Var[X] + (1-\alpha)^2 \Var[Y]
    + 2\alpha(1-\alpha)\Cov[X, Y].
    \label{eq:half_hedge_variance}
  \end{equation}
  At $\alpha = \tfrac{1}{2}$ this is
  $\tfrac{1}{4}(\Var[X] + \Var[Y] + 2\Cov[X, Y])$, smaller than
  $\Var[X]$ whenever $\Var[Y]$ and $\Cov[X, Y]$ are not too large.
\end{itemize}

A Kelly-style trade-off: give up PnL for lower variance because
the marginal seashell of PnL is worth less than the marginal
seashell of downside protection. $\alpha = \tfrac{1}{2}$ (not
Kelly-optimal) because it is robust to misestimation of $\E[X]$
and $\Var[X]$, highly uncertain in one round. Round~$5$ also
tightened the zero-crossing close for the same reason.

\paragraph{Why half, not zero?} $\alpha = 0$ removes all
constituent-direction risk, but the legs eat trading costs and
position limits, reducing throughput. Half-hedge is the compromise.

\paragraph{Why not dynamic?} Variance-minimising $\alpha$ solves a
quadratic in live $\Var[X]$, $\Var[Y]$, $\Cov[X, Y]$; the
round-sample is too small to fit. A static $\alpha$ is robust,
following the same discipline as the threshold choice.

\section{Avellaneda--Lee PCA stat-arb}
\label{sec:ch11_avellaneda_lee}

The Hedgehogs' strategy is a special case. Two assets: residual is
the spread (\cref{ch:pairs_cointegration}). $n+1$ assets with a
wrapper: residual is the basis (\cref{def:basket_basis}). The
natural extension, the one that took academic stat-arb from
pair trades to a thousand-name portfolio, is a large
constituent universe with no single wrapper, where the residual
is computed by first removing a few common factors from data.

\citet{avellaneda2010stat} formalised this in their 2010 paper
\emph{Statistical arbitrage in the U.S.\ equities market}
(\citeAL). The construction:

\begin{enumerate}
  \item Take the constituent daily-return panel
  $R \in \R^{T\times N}$ for $N \approx 500$ S\&P names.
  \item PCA on $R^\top R$; retain the first $k$ eigenvectors
  ($k$ a small single digit: $k = 1$ for the market factor,
  $k = 3$--$5$ for sector factors).
  \item For each stock $i$, regress returns on the
  principal-component series:
  $\epsilon_{i,t} = R_{i,t} - \sum_{j=1}^k \beta_{i,j} F_{j,t}$ is
  the part unexplained by common factors.
  \item Cumulate: $X_{i,t} = \sum_{s\le t} \epsilon_{i,s}$.
  Under Avellaneda--Lee each $X_{i,t}$ is approximately OU around
  zero.
  \item Apply the \cref{ch:mean_reversion_ou} $z$-score rule to
  each $X_{i,t}$. Each name is its own basket trade, with the
  factor portfolio as basket and stock $i$ as wrapper.
\end{enumerate}

The framework generalises both: pairs trading (residual from many
constituents rather than one partner) and the basket basis trade
(the wrapper is extracted statistically, not declared). We do not
derive the PCA estimator here; see the paper and \citeChan{}
(Chapter $5$, on ETF-vs-constituent strategies);
\cref{ch:ml_trading} returns to the broader ML view of factor
residuals. The point: the whole chapter so far is the special
case of Avellaneda--Lee with one factor (the wrapper), declared
loadings, and the basis as residual.

\begin{historybox}
Common-factor residuals as a trading signal predate
\citet{avellaneda2010stat} by two decades. Programme trading desks
at Salomon Brothers, Morgan Stanley, and D.E.\ Shaw in the late
1980s / early 1990s ran ``equity statistical arbitrage'' on
residuals against single-factor (CAPM) and later multi-factor
(Fama--French, BARRA) regressions, with the same OU-residual
machinery the literature later formalised.
\citet{avellaneda2010stat} is the first public-domain account,
collecting $1995$--$2008$ results and showing the strategy was
profitable until the August $2007$ ``quant quake,'' when stat-arb
funds simultaneously deleveraged and residuals decoupled from
pre-2007 OU behaviour for several days, the canonical reminder
that ``statistically stationary'' is an empirical property, not a
guarantee.
\end{historybox}

\begin{gsbox}
On a sell-side ETF arbitrage desk at Goldman Sachs, the
creation / redemption loop of \cref{sec:ch11_creation_redemption}
is the core business. The desk holds constituent inventory and
issuer connectivity (typically a daily cut-off window for AP
creations against the next NAV strike). The Strat layer owns:
\begin{itemize}
  \item \textbf{Real-time basis calculation.} A service reading
  constituent quotes, FX, dividend accruals, and borrow cost on
  any short-leg redemption hedge. Output: a live
  \emph{theoretical NAV}; its difference from ETF market price is
  the basis.
  \item \textbf{Creation / redemption cost model.} A function
  $C(\text{size},\,\text{side},\,\text{time of day})$ accounting
  for settlement lag (T+$2$ on US equities through $2024$, T+$1$
  from $2025$), borrow, and constituent-leg impact. Basis above
  cost plus margin fires the AP engine.
  \item \textbf{The automatic AP engine.} State machine that, on
  creation trigger, sends constituent buys, instructs back-office
  to create units, and queues resulting shares for market-making
  sales. Redemption mirrors. Autonomous with a one-line override.
\end{itemize}
The trader decides \emph{when to suppress}. Two canonical cases:
around the NAV-striking window (fast constituents, stale leg),
and on a name with dried-up borrow (shorting the redemption leg
is expensive). The trader pulls via an exception report.
Otherwise the engine runs the flow: SPY alone generates hundreds
of round-trips per day across the AP community; PnL is a fraction
of a basis point on each multiplied by very large notional. ETF
arb is among the most systematically run and reliably profitable
desks at a large bank. Edge has compressed as APs have
proliferated, but the fee floor remains and inventory turnover is
enormous.
\end{gsbox}

\section{A short reference implementation}
\label{sec:ch11_code}

The whole basket basis strategy is one short pure function. Input:
wall-mid prices for basket and three constituents, running mean
and std of the basis, and Olivia's \texttt{CROISSANTS} signal.
Output: one of \texttt{short\_basket}, \texttt{long\_basket},
\texttt{close}, \texttt{hold}. Drop-in for
\texttt{Trader.run(state)}. \texttt{olivia\_position} is Olivia's
inferred \texttt{CROISSANTS} position (positive/negative/zero);
\texttt{olivia\_bias} is offset $A$.

\begin{lstlisting}[caption={Frankfurt-style basket spread signal with Olivia bias.}]
def picnic_basket_spread(basket_price, croissant_price, jam_price, djembe_price,
                         running_mean, running_std, olivia_position,
                         z_entry=0.7, z_exit=0.2, olivia_bias=15):
    """Basket spread + z-score + Olivia threshold bias. Returns one action."""
    synthetic = 6 * croissant_price + 3 * jam_price + djembe_price
    basis     = basket_price - synthetic
    # Olivia adjustment: shift entry thresholds based on her croissant side.
    direction = (1 if olivia_position > 0
                 else -1 if olivia_position < 0
                 else 0)
    bias = olivia_bias * direction
    z = (basis - running_mean - bias) / max(running_std, 1e-6)
    if z >  z_entry:  return 'short_basket'   # basis is rich
    if z < -z_entry:  return 'long_basket'    # basis is cheap
    if abs(z) < z_exit: return 'close'
    return 'hold'
\end{lstlisting}

Three production notes. First, \texttt{running\_mean} and
\texttt{running\_std} are stateful and must be updated each tick
via \eqref{eq:running_mean} or an EMA; forgetting is the
\cref{pitfall:basket_premium} failure. Second, production uses
absolute-seashell thresholds, not $z$-scores ($T^{B1} = 80$,
$T^{B2} = 50$, $A = 90$ on lines $59$ and $65$);
$z_{\text{entry}} = 0.7$ above is the order-of-magnitude
equivalent divided by stationary std. Third, the returned label must be translated
by the caller into \texttt{Order} objects (typically sell the
basket's bid wall sized to the position-limit remaining; if
half-hedge \cref{sec:ch11_hh_half_hedge} is on, simultaneously
buy $\tfrac{1}{2}w_i$ of each constituent at ask).

\section{Key facts to memorise}
\label{sec:ch11_keyfacts}

\begin{itemize}
  \item \textbf{Synthetic value.} Basket $B$, weights
  $(w_1, \ldots, w_n)$ on $C_1, \ldots, C_n$: synthetic value
  $B^{\text{synth}} = \sum_i w_i C_i$, basis
  $\text{Basis}_t = B_t - B^{\text{synth}}_t$. A linear
  combination, nothing more.
  \item \textbf{Creation/redemption is the floor.} On a liquid
  AP-served ETF the basis is bounded by round-trip fee plus
  constituent liquidity costs (single-digit bps on SPY). One
  of the safest stat-arb variants because the floor is hard.
  \item \textbf{No AP $\Rightarrow$ no floor.} A basket without
  creation/redemption (proprietary indices, all Prosperity
  baskets) has no hard floor: pure cointegration play with the
  LTCM-box caveats of \cref{ch:pairs_cointegration}.
  \item \textbf{Cointegration vector is declared, not fitted.}
  With published weights, $\bm\beta = (1, -w_1, \ldots, -w_n)$
  is known: no Engle--Granger/Johansen. Walk-forward refit
  applies to running mean and variance, not to the slope.
  \item \textbf{Static vs dynamic hedge.} Static = published
  weights (AP or simulator enforces). Dynamic = fitted weights
  (loose coupling, drifting effective ratio).
  \item \textbf{Trading rule.} Short basis at
  $z_t > +z_{\text{entry}}$, long at $z_t < -z_{\text{entry}}$,
  close at $|z_t| < z_{\text{exit}}$: the same OU $z$-score
  machinery as \cref{ch:mean_reversion_ou}.
  \item \textbf{Frankfurt PICNIC compositions.}
  \texttt{PICNIC\_BASKET1} $= 6$\,\texttt{CROISSANTS} $+
  3$\,\texttt{JAMS} $+ 1$\,\texttt{DJEMBE};
  \texttt{PICNIC\_BASKET2} $= 4$\,\texttt{CROISSANTS} $+
  2$\,\texttt{JAMS}. \texttt{DJEMBES} only in basket~$1$. Source:
  \texttt{ETF\_CONSTITUENT\_FACTORS = [[6, 3, 1], [4, 2, 0]]},
  line $57$ of \texttt{FrankfurtHedgehogs\_polished.py}.
  \item \textbf{Frankfurt thresholds.}
  \texttt{BASKET\_THRESHOLDS = [80, 50]} (raw seashells) and
  Olivia adjustment \texttt{ETF\_THR\_INFORMED\_ADJS = [90, 90]};
  both thresholds shift by $d_t \cdot A$.
  \item \textbf{Premium drift.} Basis mean is non-zero and
  slowly drifting. Fix: online running-mean update
  \eqref{eq:running_mean} before thresholding; else the rule is
  systematically biased.
  \item \textbf{Olivia cross-product.} The same bot drives
  \texttt{CROISSANTS} and the basket basis. Shifting the band
  by $\pm A$ was worth a material fraction of basket PnL.
  \item \textbf{Headline PnL.} $40$k--$60$k seashells/round on
  baskets, plus $20$k/round on \texttt{CROISSANTS} alone.
  \item \textbf{Round 5 half-hedge.} Convex combination of
  unhedged and fully-hedged strategies: linear in return,
  sub-linear in variance via \eqref{eq:half_hedge_variance}.
  Kelly-style trade-off (\cref{ch:kelly_risk}).
  \item \textbf{Avellaneda--Lee.} Basket basis is the one-factor
  special case of PCA stat-arb: with $k$ PCA factors removed on
  the constituent return panel, every stock has its own residual
  basis and trades on the same OU $z$-score rule
  (\cref{ch:ml_trading}).
\end{itemize}
